\documentclass[11pt,a4paper]{article}
\usepackage[italian]{babel}
\usepackage[margin=2.5cm]{geometry}
\usepackage{parskip}
\usepackage{amsmath, amssymb, amsthm}
\usepackage{booktabs}
\usepackage{array}
\usepackage{xcolor}
\usepackage{needspace}
\usepackage{enumitem}
\usepackage{listings}
\usepackage{tikz}
\usepackage[most]{tcolorbox}
\usepackage[hidelinks]{hyperref}
\usetikzlibrary{decorations.pathreplacing, shapes.geometric, arrows.meta, positioning}

% Niente vedove/orfane: i paragrafi non lasciano righe isolate a fine/inizio pagina
\widowpenalty=10000
\clubpenalty=10000

% Elenchi più compatti
\setlist{itemsep=2pt, parsep=0pt, topsep=4pt, partopsep=0pt}

% --- Riquadri con bordo nero, senza numero (si spezzano tra due pagine) ---
% Uso: \begin{definizione} ... \end{definizione}
%      \begin{teorema}[Nome] ... \end{teorema}  (il nome è facoltativo)
\tcbset{nota/.style={enhanced jigsaw, breakable, colback=white,
  colframe=black, boxrule=0.5pt, sharp corners}}
\NewTColorBox{definizione}{}{nota,
  before upper={\textbf{Definizione.}\ }}
\NewTColorBox{teorema}{o}{nota,
  before upper={\textbf{Teorema\IfValueT{#1}{ (#1)}.}\ }}
\NewTColorBox{proposizione}{o}{nota,
  before upper={\textbf{Proposizione\IfValueT{#1}{ (#1)}.}\ }}
\NewTColorBox{proprieta}{o}{nota,
  before upper={\textbf{Proprietà\IfValueT{#1}{ (#1)}.}\ }}
\NewTColorBox{assioma}{o}{nota, boxrule=1pt,
  before upper={\textbf{Assioma\IfValueT{#1}{ (#1)}.}\ }}

% --- Ambienti semplici (senza riquadro) ---
% Il titolo sta in cima al blocco, anche se il contenuto inizia con un riquadro o
% con codice; e non resta mai da solo in fondo a una pagina.
% Uso: \begin{esempio}[Titolo facoltativo] ... \end{esempio}
\newcommand{\newrunin}[2]{%
  \NewDocumentEnvironment{#1}{o}{%
    \par\Needspace{4\baselineskip}\addvspace{\medskipamount}\noindent
    \textbf{#2}\IfValueT{##1}{ (##1)}\textbf{.}\ \ignorespaces}%
  {\par\addvspace{\medskipamount}}}
\newrunin{esempio}{Esempio}
\newrunin{esercizio}{Esercizio}
\NewTColorBox{osservazione}{}{nota,
  before upper={\textbf{Osservazione.}\ }}

% V e F nelle tavole di verità
\newcommand{\V}{\textsf{V}}
\newcommand{\F}{\textsf{F}}

% --- Codice C ---
% Uso: \begin{codice} ... \end{codice}      codice C numerato
%      \begin{comando} ... \end{comando}    comandi da terminale
%      \begin{risultato} ... \end{risultato} output di un programma
%      \lstinline|printf("ciao");|           codice dentro al testo
\lstdefinestyle{c}{
  language=C,
  basicstyle=\ttfamily\small,
  keywordstyle=\color{blue}\bfseries,
  commentstyle=\color{gray}\itshape,
  stringstyle=\color{red!70!black},
  numbers=left,
  numberstyle=\tiny\color{gray},
  numbersep=8pt,
  columns=flexible,
  keepspaces=true,
  breaklines=true,
  showstringspaces=false,
  tabsize=4,
  morekeywords={include, define},
  % lettere accentate nei commenti
  literate={à}{{\`a}}1 {è}{{\`e}}1 {é}{{\'e}}1 {ì}{{\`i}}1
           {ò}{{\`o}}1 {ù}{{\`u}}1 {È}{{\`E}}1 {À}{{\`A}}1
}
\lstset{style=c}
\newtcblisting{codice}{listing only, nota, left=6mm,
  listing options={style=c}}
\newtcblisting{comando}{listing only, nota,
  listing options={basicstyle=\ttfamily\small, numbers=none, language={},
    columns=flexible, keepspaces=true, breaklines=true}}
\newtcblisting{risultato}{listing only, enhanced jigsaw, breakable,
  colback=black!5, colframe=black, boxrule=0.5pt, sharp corners,
  listing options={basicstyle=\ttfamily\small, numbers=none, language={},
    columns=flexible, keepspaces=true, breaklines=true}}

% --- Diagrammi di flusso (simboli standard) ---
\tikzset{
  terminale/.style = {ellipse, draw, minimum width=2.5cm, minimum height=0.9cm},
  io/.style        = {trapezium, trapezium left angle=70, trapezium right angle=110,
                      draw, minimum width=2.5cm, minimum height=0.9cm},
  processo/.style  = {rectangle, draw, minimum width=2.5cm, minimum height=0.9cm},
  decisione/.style = {diamond, draw, aspect=2, minimum width=2.5cm},
  freccia/.style   = {thick, -{Stealth}}
}

\title{Algoritmi e diagrammi di flusso}
\author{Marco Cerbella}
\date{Lezione 1 -- 22 settembre 2026}

\begin{document}
\maketitle

\section{Algoritmi e diagrammi di flusso}

\begin{definizione}
Un \emph{algoritmo} è una procedura che risolve un problema attraverso una
sequenza finita di passi elementari e non ambigui. Spesso comporta la
ripetizione di alcune operazioni.
\end{definizione}

\begin{esempio}
Il primo algoritmo celebre è quello di Euclide (\emph{Elementi}, 300 a.C.) per
il massimo comun divisore.
\end{esempio}

\begin{definizione}
Un \emph{diagramma di flusso} (\emph{flowchart}) è la rappresentazione grafica
di un algoritmo. Le operazioni si leggono dall'alto verso il basso, seguendo le
frecce. I simboli sono standardizzati (ANSI/ISO).
\end{definizione}

I simboli principali sono:
\begin{itemize}
    \item \textbf{ellisse}: inizio e fine;
    \item \textbf{parallelogramma}: input e output (I/O);
    \item \textbf{rettangolo}: operazione (\emph{processo});
    \item \textbf{rombo}: decisione (condizione vero/falso).
\end{itemize}

\begin{center}
\begin{tikzpicture}[node distance=0.8cm]
  \node (inizio) [terminale] {Inizio};
  \node (input)  [io, below=of inizio] {Input};
  \node (op)     [processo, below=of input] {Operazioni};
  \node (output) [io, below=of op] {Output};
  \node (fine)   [terminale, below=of output] {Fine};

  \draw [freccia] (inizio) -- (input);
  \draw [freccia] (input) -- (op);
  \draw [freccia] (op) -- (output);
  \draw [freccia] (output) -- (fine);
\end{tikzpicture}
\end{center}

Lo stesso schema in linguaggio \textit{C}:
\begin{codice}
// Importare le librerie necessarie
#include <stdio.h>

// Funzione principale
int main() {
    // Lettura dell'input e operazioni

    // Stampa del risultato
    return 0;
}
\end{codice}

\section{Programmare}

\begin{definizione}
\emph{Programmare} significa tradurre un algoritmo in una notazione, un
linguaggio di programmazione, che un elaboratore può eseguire. Un
\emph{programma} è un insieme di istruzioni.
\end{definizione}

I passi sono: partire dal problema, \emph{pensare all'algoritmo}, scegliere il
linguaggio giusto e solo dopo programmare.

\section{Esempio: massimo comun divisore}

Algoritmo di Euclide (versione con sottrazioni): finché $b \neq 0$, si
sostituisce il maggiore tra $a$ e $b$ con la differenza tra i due. Alla fine $a$
contiene il MCD.

\begin{center}
\begin{tikzpicture}
  \node (s) at (0,0) [terminale] {Inizio};
  \node (in) at (0,-1.3) [io] {Leggi $a$, $b$};
  \node (d1) at (0,-3.3) [decisione] {$b \neq 0$};
  \node (d2) at (0,-5.8) [decisione] {$a > b$};
  \node (pa) at (-2.3,-8) [processo] {$a = a - b$};
  \node (pb) at (2.3,-8) [processo] {$b = b - a$};
  \node (j) at (0,-9.5) {};
  \node (m) at (0,-2.3) {};
  \node (out) at (4.3,-3.3) [io] {Stampa $a$};
  \node (e) at (4.3,-4.8) [terminale] {Fine};

  \draw [freccia] (s) -- (in);
  \draw [freccia] (in) -- (d1);
  \draw [freccia] (d1) -- node[right]{Sì} (d2);
  \draw [freccia] (d1.east) -- node[above]{No} (out.west);
  \draw [freccia] (out) -- (e);
  \draw [freccia] (d2.west) -| node[above, pos=0.2]{Sì} (pa.north);
  \draw [freccia] (d2.east) -| node[above, pos=0.2]{No} (pb.north);
  \draw (pa.south) |- (j.center);
  \draw (pb.south) |- (j.center);
  \draw [freccia] (j.center) -- (-4.4,-9.5) -- (-4.4,-2.3) -- (m.center);
\end{tikzpicture}
\end{center}

E il programma C corrispondente:
\begin{codice}
#include <stdio.h>

int main() {
    int a, b;

    printf("Enter first positive integer: \n");
    scanf("%d", &a);
    printf("Enter second positive integer: \n");
    scanf("%d", &b);

    while (b != 0) {
        if (a > b)
            a = a - b;
        else
            b = b - a;
    }
    printf("GCD = %d\n", a);

    return 0;
}
\end{codice}

\begin{osservazione}
Il ciclo \lstinline|while| corrisponde al rombo $b \neq 0$ con la freccia che
torna indietro; l'\lstinline|if|/\lstinline|else| corrisponde al secondo rombo.
\end{osservazione}
\end{document}
